\section{A transform-level Laguerre attack}\label{sec:transformattack}
The Riemann hypothesis asks about zeros of the completed
xi-function, rather than the curvature of its kernel.
With the present Csordas--Planat--Sol\'e normalization,
write $K(u)=\Phi(u)$ and
\begin{equation}\label{eq:fourier-actual}
F(z)=\int_{\mathbb R}K(u)e^{izu}\,du
     =\frac14\,\xi\!\left(\frac12+\frac{iz}{2}\right).
\end{equation}
The factor $1/4$ and the half-frequency must be retained:
the first real zero of $F$ is near $z=28.269$, not $14.134$.
The scaling does not affect signs of generalized Laguerre
expressions.

For a sufficiently rapidly decreasing positive even kernel,
set
\begin{equation}
L_n[F](x)=\frac{(-1)^n}{(2n)!}
\sum_{j=0}^{2n}(-1)^j {2n\choose j}
F^{(j)}(x)F^{(2n-j)}(x).
\end{equation}
Differentiating the Fourier integral and symmetrizing
yields an exact correlation representation:
\begin{equation}\label{eq:all-order-corr}
L_n[F](x)=\frac{1}{(2n)!}
\iint_{\mathbb R^2}(u-v)^{2n}K(u)K(v)
\cos(x(u+v))\,du\,dv.
\end{equation}

\begin{proposition}[Central positivity and a local band]
Let $K$ be real, even, positive almost everywhere and have
the moments required in \eqref{eq:all-order-corr}.
Define the positive quantities
\begin{align}
A_n&=\frac{1}{(2n)!}\iint (u-v)^{2n}K(u)K(v)\,du\,dv,\\
B_n&=\frac{1}{(2n)!}\iint (u-v)^{2n}(u+v)^2K(u)K(v)\,du\,dv.
\end{align}
Then $L_n[F](0)=A_n>0$ and for every real $x$,
\begin{equation}\label{eq:local-positive}
L_n[F](x)\ge A_n-\frac{x^2}{2}B_n.
\end{equation}
In particular $L_n[F](x)>0$ for
$|x|<\sqrt{2A_n/B_n}$ whenever $B_n>0$.
\end{proposition}
\begin{proof}
Equation~\eqref{eq:all-order-corr} is the coefficient of
$y^{2n}$ in
$F(x+iy)F(x-iy)$, using symmetry under $u\leftrightarrow v$.
The strict inequality at $x=0$ follows from positivity of
the integrand away from the measure-zero diagonal.
Finally $\cos\theta\ge 1-\theta^2/2$ for every real $\theta$,
and integration against the nonnegative correlation weight
gives~\eqref{eq:local-positive}.
\end{proof}
This observation is a diagnostic obstruction, not an RH
proof: central positivity holds even for the explicit
counterexamples of the preceding section. The all-order
correlation construction has prior art in
Dimitrov--Xu~\cite{DX}.

For example, with $M_j=\int_{\mathbb R}u^j K(u)\,du$
and $M_{2j+1}=0$, direct expansion at order one gives
\begin{equation}\label{eq:n1moments}
A_1=M_0M_2,\qquad B_1=M_0M_4-M_2^2.
\end{equation}
The latter is nonnegative by Cauchy--Schwarz.

\section{Numerical exploration of the genuine transform}
The reproducible experiment
\texttt{certificates/generalized\_laguerre\_scan.py}
uses 65-decimal-digit \texttt{mpmath} quadrature for the
Jacobi-theta representation. It evaluates
$L_n[F](x)$ for $n=1,\ldots,5$ at
$x\in\{0,5,10,15,20,25,28,30\}$.
All 40 values were \emph{numerically} positive.
At the endpoints, $L_1[F](0)\simeq
1.78434856788\times 10^{-4}$ and
$L_1[F](30)\simeq 1.09904761087\times10^{-8}$.
Independent differentiation of the completed zeta function
in \eqref{eq:fourier-actual} at $x=0,10,28$ agrees to roughly
$10^{-60}$ relative precision for derivatives through order four.

The numerically estimated local positivity radii
$\sqrt{2A_n/B_n}$ for $n=1,\ldots,8$ are, respectively,
\[
9.8314,\ 10.3132,\ 10.7562,\ 11.1665,\
11.5491,\ 11.9081,\ 12.2466,\ 12.5674.
\]
These are exploratory evaluations, \emph{not} rigorously
enclosed moment integrals or proofs of positivity outside
the sampled points. Their growth suggests looking for
an $n$-uniform, noncentral correlation estimate, but no
such estimate is claimed here.

\section{Lean expansion and the remaining analytic bridge}
The new module \texttt{XiKernel/TransformMoment.lean}
proves a finite-sum central-moment positivity theorem for
every natural number $n$, the elementary quadratic gate,
and exact order-one moment identities. Its hypotheses
involve abstract strictly positive even moments;
the analytic identification with the real theta-kernel
integrals is \emph{not} formalized.

The module \texttt{XiKernel/RationalJet.lean}
defines a finite Taylor-jet recurrence and computes the
fifth iterate of rational 12-decimal approximants to the
eleven genuine theta Taylor coefficients. GitHub Lean CI
passed, but its \texttt{native\_decide} proof depends on
\texttt{Lean.ofReduceBool} and \texttt{Lean.trustCompiler},
as confirmed by an explicit \texttt{\#print axioms} audit.
It is compiler-assisted reflection, \emph{not} a small-kernel-only
certificate. This establishes only a rational sign for the
\emph{approximate} jet; it cannot be substituted for
Theorem~\ref{thm:five} without a formally checked
analytic coefficient enclosure and stability theorem.

To certify Theorem~\ref{thm:five} \emph{end to end in Lean},
three missing links must all be discharged:
(i) the theta-series differentiation and Taylor-coefficient
identities within Mathlib, including modular evenness;
(ii) formal rational enclosures for $\pi$, exponential
truncation, and every omitted $m\ge13$ summand;
(iii) sound interval propagation from these enclosures
to the actual fifth iterate and the strict negative bound.
The existing standard-library Python calculation supplies
a reproducible exact computational certificate but is not
a kernel-checked Lean derivation of (i)--(iii).
No \texttt{sorry}, fabricated axiom, or numerical
observation is to be represented as closing these steps.
